\section{Problem Statement, Threat Model and Design Goals}
\label{sec:threat}

\subsection{Problem statement}
A university runs an OIDC identity provider for its students and a set of services of different trustworthiness. We want each service to receive only the attributes it is entitled to and has declared a purpose for, to be unable to link a student across services by identifier, and to be held to its declared purpose and retention period, as far as the institute can observe. Every release, and every observable attempt to break these rules, should leave a record that the student can inspect and that nobody, including the operator of the system, can later change without detection. A service that breaks the rules should automatically receive less in future. Services should be able to keep using standard OIDC.

\subsection{Parties}
\begin{itemize}
\item \emph{Students} log in through the IdP and want to know what each service received.
\item The \emph{IdP} (Keycloak) authenticates students, holds their attribute records, and issues tokens.
\item \emph{TrustTrail} verifies declarations, releases attribute values, keeps the ledger, computes trust and anchors the ledger.
\item \emph{Services} are relying parties of differing trust: the library, the LMS and the examination portal in our deployment, and possibly third-party services.
\item The \emph{operator} administers TrustTrail and its database, and can reset a service's trust.
\item A \emph{verifier} is anyone who checks a receipt or the ledger against the chain: a student, an auditor or the institute's grievance officer.
\item The \emph{chain} is an Ethereum network on which the anchor contract is deployed.
\end{itemize}

\subsection{Adversaries}
We consider the following adversaries. Each is a role that a party may take, not necessarily a separate attacker.

\begin{enumerate}
\item[\textbf{A1}] \emph{Over-asking service.} A service requests attributes that its declaration does not cover, or a profile above the one its agreement allows.
\item[\textbf{A2}] \emph{Retaining service.} A service keeps using a declaration after its promised deletion date, for example by requesting the data again.
\item[\textbf{A3}] \emph{Colluding services.} Two or more services share what they received in order to link a student's activity across them, using identifiers, session values or authentication timestamps in the tokens.
\item[\textbf{A4}] \emph{Repudiating service.} A service later denies having promised a purpose or a deletion date, or denies having received data.
\item[\textbf{A5}] \emph{Insider with database access.} Someone with write access to TrustTrail's database file, possibly the operator, edits, deletes or reorders ledger entries, rewrites the whole chain from some entry onward, or cuts off the newest entries. We assume this adversary can remove the database's append-only triggers.
\item[\textbf{A6}] \emph{Forger.} An attacker presents a forged access token to obtain attributes, or forges a receipt to claim a release that did not happen.
\end{enumerate}

\subsection{Assumptions}
\begin{itemize}
\item The IdP is honest and correctly configured: it authenticates students correctly, signs tokens with keys that only it holds, and grants a client only the scopes assigned to it.
\item TrustTrail runs its code faithfully \emph{at the time of a release}. Adversary A5 may tamper with its stored data afterwards, but we do not consider a TrustTrail that lies while writing an entry; Section~\ref{sec:discussion} discusses this case.
\item Each service keeps its declaration-signing key and its client secret private. A stolen key lets the thief sign in the service's name. Registering a key requires only the service's client secret, and the key registry is held by TrustTrail, so the client secret is the root of trust for non-repudiation.
\item TrustTrail holds the receipt-signing key and the anchoring key, and owns the anchor contract. The operator can therefore anchor whatever its ledger contains; what it cannot do is change or remove a batch once it is on chain.
\item The anchor contract runs on a chain that the operator cannot rewrite. A local development chain does not meet this assumption; a public network does.
\item Communication uses TLS in a real deployment. Our local deployment uses plain HTTP between containers on one machine.
\item Hash functions (SHA-256) and signature schemes (ECDSA P-256, RSA) are secure.
\end{itemize}

\subsection{Design goals}
Table~\ref{tab:goals} lists the design goals, the adversaries each goal addresses, and the survey characteristics~\cite{werner2017survey} each goal corresponds to.

\begin{table}[t]
\centering
\caption{Design goals}
\label{tab:goals}
\small
\begin{tabularx}{\linewidth}{@{}l >{\raggedright\arraybackslash}X >{\raggedright\arraybackslash}p{1.8cm} >{\raggedright\arraybackslash}p{3.6cm}@{}}
\toprule
Goal & Statement & Adversaries & Survey characteristic \\
\midrule
G1 & A service receives at most the attributes of the profile allowed by its current tier, and only those its declaration covers. & A1 & Minimization, controllability \\
G2 & No identifier or session value in the tokens or releases lets services link a student across services. (Released attribute values may still link; see non-goals.) & A3 & Obfuscation (pseudonym), unlinkability \\
G3 & Every release is covered by a service-signed declaration of purpose and deletion date; requests outside the declaration or after the date are refused. & A1, A2 & Use limitation \\
G4 & A service cannot deny a declaration signed with its registered key, and TrustTrail cannot deny a release it receipted. & A4 & Accountability (policies) \\
G5 & Any change, deletion or truncation of entries that have been anchored on a chain the operator cannot rewrite is detectable, also when the operator makes it; newer entries are protected by the chain and by receipts. & A5 & Accountability (traceability) \\
G6 & Observed violations lower a service's trust and, through its tier, what it can obtain. & A1, A2 & Trust \\
G7 & A student can see every release and refused rule break that concerns them, and verify each one independently. & A4, A5 & Transparency, user-friendliness \\
G8 & Services use standard OIDC plus signed JSON; the added latency per login stays in the order of milliseconds. & -- & (deployability) \\
\bottomrule
\end{tabularx}
\end{table}

\subsection{Non-goals}
Some threats are outside TrustTrail's scope. We list them explicitly so the guarantees are not over-read.
\begin{itemize}
\item \emph{Use after release.} Once a service holds data, TrustTrail cannot see whether it is used for the declared purpose or deleted on time. Rule~2 catches only a service that asks again.
\item \emph{IdP tracking.} Keycloak sees every login. The protocols in Section~\ref{sec:related} address this.
\item \emph{Traffic analysis.} Services that compare login times or network addresses may still correlate pseudonyms. TrustTrail removes identifiers, not timing.
\item \emph{Linking through attribute values.} The attributes that a profile releases, such as department, programme and year, can together single out a student in a small cohort, and the total profile releases name, e-mail address and roll number. Colluding services that hold such values can join their records without any identifier. TrustTrail limits which attributes each service receives; it does not make the released values unlinkable.
\item \emph{A dishonest release service.} If TrustTrail itself omits a release from the ledger, anchoring cannot reveal an entry that was never written; only the receipt held by the student is counter-evidence.
\item \emph{Availability.} Denial-of-service against Keycloak, TrustTrail or the chain is not addressed.
\end{itemize}
